\section{Introduction}

Multimodal models increasingly interpret people from several affective channels at
once. In favorable cases, facial behavior, language, and voice reinforce one another.
In harder cases, a positive utterance may accompany visible strain, or a calm face may
co-occur with distressed speech. Such conflict matters because a model can produce a
fluent response while grounding its affective interpretation in the wrong cue. An
output-level label tells us whether this affective misread occurred, but not what
internal configuration preceded it. Multimodal affect datasets such as CH-SIMS v2
provide a basis for studying how these cues interact \cite{liu2022chsims}; our concern
is the organization of that interaction inside a generative model before it speaks.
Recent conflict-centered benchmarks show systematic modality bias and performance
degradation under affective conflict \cite{han2025benchmarking,sun2026emomm}, but
output accuracy and attention summaries do not by themselves describe the geometry
of the complete pre-generation relation.

We therefore ask a two-part question: how is multimodal affective conflict organized
in the model's internal state before the first generated token, and how are these
pre-generation state structures associated with subsequent affective misreads? This
question differs from moving a binary error classifier to an earlier time step. It
also differs from learning another Conflict detector. Conflict/Aligned describes a
property of the registered multimodal input, whereas Misread/Non-misread describes a
model's later behavior relative to a reference affect description. We use the former
to learn a stable relation space and reserve the latter for an independent behavioral
analysis. Misread labels are never used to train the state encoder.

Our measurement point is \(t_0\), the last conditioning state before the first
generated token. For every sample, we construct two unimodal conditions \(M_1\) and
\(M_2\), together with the joint condition \(M_{12}\), and extract their complete
layer trajectories under \(P\) registered, semantically equivalent prompts. The
Trajectory Manifold Encoder (TME) maps each condition trajectory to a unit-sphere
embedding and maps the ordered three-condition relation to a sample-level relation
embedding. TME learns this relation from sample-level Conflict/Aligned supervision;
the class label is not copied onto any individual condition trajectory. After
training, the encoder is frozen. Its condition-level embeddings and sample-level
relation embeddings are cached separately because they answer different questions.

The condition-level embeddings define three complementary State Indices. State
Dispersion \(S\) asks whether the internal state remains stable when the same task is
phrased in equivalent ways. Modality Split \(D\) asks whether the two unimodal inputs
pull the model toward distinct internal directions. Signed Joint Lean \(R\) asks
which unimodal direction the joint input approaches; by convention, \(R>0\) denotes
V-lean and \(R<0\) denotes T/A-lean. Put simply, \(S\) asks whether the model stands
steadily, \(D\) asks whether the modalities pull apart, and \(R\) asks which way the
joint state leans. They are coordinates of different phenomena, not three versions of
one severity variable, and we do not combine them into a universal Misread
probability.

Aligned calibration data partition these coordinates into four State Patterns.
Confusion means that equivalent prompts move the state too much to assign a stable
direction. Among stable states, Consensus means that the unimodal states do not
separate significantly; Balanced means that they separate but the joint state does
not lean significantly toward either one; and Dominant means that a significant
split is accompanied by a significant directional lean. State Pattern is therefore
a mechanism-motivated descriptive partition of pre-generation configurations, not a
taxonomy of errors. A Confusion state may still yield a correct description, and a Dominant state
may lean toward either the reference-supported cue or a misleading surface cue.
Likewise, Consensus may reflect genuinely aligned evidence or a failure to encode an
objective conflict. The patterns describe how a model organizes evidence, not whether
its eventual answer must be right or wrong.

This distinction determines our evidence chain. First, we measure whether Conflict
inputs have a different prevalence of affective misreads from Aligned inputs. Second,
we test whether Conflict and Aligned inputs differ in their State Index and State
Pattern distributions. Third, within Conflict samples only, we measure how \(S\),
\(D\), State Pattern, and the direction of \(R\) are empirically associated with
Misread outcomes. Signed Joint Lean primarily explains direction; its absolute value
is not treated as a general error score. Fourth, we freeze each learned representation
and train the same controlled Conflict-only Misread probe to test how much behavioral
information the representation retains. Finally, we evaluate whether state-guided
responses reduce response-level Affective Misattribution by preserving overlooked
cues. This last step does not claim to change the base model's earlier Diagnostic
Misread.

Our representation comparison follows the same boundary. Single-Point,
Trajectory MLP, and TME use the same master split, Conflict/Aligned data, Conflict
supervision budget, three conditions, probe architecture, probe budget, and held-out
intersection. Each method nevertheless retains its native representation and training
objective: the two baselines use ordinary classifiers, whereas TME uses its spherical
ordered-relation module and Proxy Anchor objective. The comparison therefore concerns
each complete method. Without component ablations, an advantage cannot be assigned
solely to full-layer trajectories, the manifold geometry, or Proxy Anchor.

The contributions of this work are as follows:
\begin{itemize}
  \item We formulate affective conflict analysis as the measurement of
  pre-generation state structure and its association with subsequent Misread
  behavior, explicitly separating input-side Conflict/Aligned supervision from
  output-side Misread evaluation.
  \item We introduce TME and a registered three-condition, multi-prompt protocol that
  exports condition-level spherical states and sample-level relation embeddings at
  \(t_0\).
  \item We define State Dispersion, Modality Split, signed Joint Lean, and four
  calibrated State Patterns that explain complementary internal phenomena without
  constructing a scalar error score or an ordered error hierarchy.
  \item We establish a matched evaluation of complete frozen representations with a
  Conflict-only Misread probe, and separate this representation analysis from the
  downstream evaluation of response-level Affective Misattribution.
\end{itemize}
